\section{Ample and very ample sheaves relative to finite type morphisms}
\label{section-ample-finite-type}

\noindent
In fact most of the material in this section is about the notion of
a (quasi-)projective morphism which we have not defined yet.

\begin{lemma}
\label{lemma-very-ample-finite-type-over-affine}
Let $f : X \to S$ be a morphism of schemes.
Let $\mathcal{L}$ be an invertible sheaf on $X$.
Assume that
\begin{enumerate}
\item the invertible sheaf $\mathcal{L}$ is very ample on $X/S$,
\item the morphism $X \to S$ is of finite type, and
\item $S$ is affine.
\end{enumerate}
Then there exist an $n \geq 0$ and an immersion
$i : X \to \mathbf{P}^n_S$ over $S$ such that
$\mathcal{L} \cong i^*\mathcal{O}_{\mathbf{P}^n_S}(1)$.
\end{lemma}

\begin{proof}
Assume (1), (2) and (3).
Condition (3) means $S = \Spec(R)$ for some ring $R$.
Condition (1) means by definition
there exists a quasi-coherent $\mathcal{O}_S$-module
$\mathcal{E}$ and an immersion $\alpha : X \to \mathbf{P}(\mathcal{E})$
such that $\mathcal{L} = \alpha^*\mathcal{O}_{\mathbf{P}(\mathcal{E})}(1)$.
Write $\mathcal{E} = \widetilde{M}$ for some $R$-module $M$.
Thus we have
$$
\mathbf{P}(\mathcal{E}) = \text{Proj}(\text{Sym}_R(M)).
$$
Since $\alpha$ is an immersion, and since the topology of
$\text{Proj}(\text{Sym}_R(M))$ is generated by the standard
opens $D_{+}(f)$, $f \in \text{Sym}_R^d(M)$, $d \geq 1$,
we can find for each $x \in X$ an
$f \in \text{Sym}_R^d(M)$, $d \geq 1$, with $\alpha(x) \in D_{+}(f)$
such that
$$
\alpha|_{\alpha^{-1}(D_{+}(f))} : \alpha^{-1}(D_{+}(f)) \to D_{+}(f)
$$
is a closed immersion.
Condition (2) implies $X$ is quasi-compact. Hence we can find
a finite collection of elements
$f_j \in \text{Sym}_R^{d_j}(M)$, $d_j \geq 1$
such that for each $f = f_j$ the displayed map above
is a closed immersion and such that $\alpha(X) \subset \bigcup D_{+}(f_j)$.
Write $U_j = \alpha^{-1}(D_{+}(f_j))$. Note that $U_j$ is affine
as a closed subscheme of the affine scheme $D_{+}(f_j)$.
Write $U_j = \Spec(A_j)$. Condition (2) also implies that
$A_j$ is of finite type over $R$, see
Lemma \ref{lemma-locally-finite-type-characterize}.
Choose finitely many $x_{j, k} \in A_j$ which
generate $A_j$ as a $R$-algebra. Since $\alpha|_{U_j}$ is a closed
immersion we see that $x_{j, k}$ is the image of an element
$$
f_{j, k}/f_j^{e_{j, k}} \in \text{Sym}_R(M)_{(f_j)}
=
\Gamma(D_{+}(f_j), \mathcal{O}_{\text{Proj}(\text{Sym}_R(M))}).
$$
Finally, choose $n \geq 1$ and elements $y_0, \ldots, y_n \in M$ such that each
of the polynomials $f_j, f_{j, k} \in \text{Sym}_R(M)$ is a polynomial
in the elements $y_t$ with coefficients in $R$.
Consider the graded ring map
$$
\psi : R[Y_0, \ldots, Y_n] \longrightarrow \text{Sym}_R(M),
\quad Y_i \longmapsto y_i.
$$
Denote $F_j$, $F_{j, k}$ the elements of $R[Y_0, \ldots, Y_n]$ such
that $\psi(F_j) = f_j$ and $\psi(F_{j, k}) = f_{j, k}$.
By Constructions, Lemma \ref{constructions-lemma-morphism-proj}
we obtain an open subscheme
$$
U(\psi) \subset \text{Proj}(\text{Sym}_R(M))
$$
and a morphism
$r_\psi : U(\psi) \to \mathbf{P}^n_R$. This morphism
satisfies $r_\psi^{-1}(D_{+}(F_j)) = D_{+}(f_j)$, and hence we see
that $\alpha(X) \subset U(\psi)$. Moreover, it is clear
that
$$
i = r_\psi \circ \alpha : X \longrightarrow \mathbf{P}^n_R
$$
is still an immersion since
$i^\sharp(F_{j, k}/F_j^{e_{j, k}}) = x_{j, k} \in
A_j = \Gamma(U_j, \mathcal{O}_X)$
by construction. Moreover, the morphism $r_\psi$ comes
equipped with a map
$\theta : r_\psi^*\mathcal{O}_{\mathbf{P}^n_R}(1)
\to \mathcal{O}_{\text{Proj}(\text{Sym}_R(M))}(1)|_{U(\psi)}$
which is an isomorphism in this case (for construction $\theta$
see lemma cited above; some details omitted).
Since the original map $\alpha$ was assumed to have the
property that
$\mathcal{L} = \alpha^*\mathcal{O}_{\text{Proj}(\text{Sym}_R(M))}(1)$
we win.
\end{proof}

\begin{lemma}
\label{lemma-quasi-affine-finite-type-over-S}
Let $\pi : X \to S$ be a morphism of schemes.
Assume that $X$ is quasi-affine and that $\pi$ is locally of finite type.
Then there exist $n \geq 0$ and an immersion $i : X \to \mathbf{A}^n_S$
over $S$.
\end{lemma}

\begin{proof}
Let $A = \Gamma(X, \mathcal{O}_X)$. By assumption $X$ is quasi-compact
and is identified with an open subscheme of $\Spec(A)$, see
Properties, Lemma \ref{properties-lemma-quasi-affine}.
Moreover, the set of opens $X_f$, for those $f \in A$ such that $X_f$ is
affine, forms a basis for the topology of $X$, see the proof of
Properties, Lemma \ref{properties-lemma-quasi-affine}.
Hence we can find a finite number of $f_j \in A$, $j = 1, \ldots, m$ such that
$X = \bigcup X_{f_j}$, and such that $\pi(X_{f_j}) \subset V_j$ for
some affine open $V_j \subset S$. By
Lemma \ref{lemma-locally-finite-type-characterize}
the ring maps $\mathcal{O}(V_j) \to \mathcal{O}(X_{f_j}) = A_{f_j}$
are of finite type. Thus we may choose $a_1, \ldots, a_N \in A$ such that
the elements $a_1, \ldots, a_N, 1/f_j$ generate
$A_{f_j}$ over $\mathcal{O}(V_j)$ for each $j$. Take $n = m + N$ and
let
$$
i : X \longrightarrow \mathbf{A}^n_S
$$
be the morphism given by the global sections
$f_1, \ldots, f_m, a_1, \ldots, a_N$ of the structure sheaf of $X$.
Let $D(x_j) \subset \mathbf{A}^n_S$ be the open subscheme where the
$j$th coordinate function is nonzero.
Then for $1 \leq j \leq m$ we have $i^{-1}(D(x_j)) = X_{f_j}$ and
the induced morphism $X_{f_j} \to D(x_j)$ factors through the affine
open $\Spec(\mathcal{O}(V_j)[x_1, \ldots, x_n, 1/x_j])$
of $D(x_j)$. Since the ring map
$\mathcal{O}(V_j)[x_1, \ldots, x_n, 1/x_j] \to A_{f_j}$ is
surjective by construction we conclude that $i^{-1}(D(x_j)) \to D(x_j)$
is an immersion as desired.
\end{proof}

\begin{lemma}
\label{lemma-quasi-projective-finite-type-over-S}
Let $f : X \to S$ be a morphism of schemes.
Let $\mathcal{L}$ be an invertible sheaf on $X$.
Assume that
\begin{enumerate}
\item the invertible sheaf $\mathcal{L}$ is ample on $X$, and
\item the morphism $X \to S$ is locally of finite type.
\end{enumerate}
Then there exists a $d_0 \geq 1$ such that for every $d \geq d_0$
there exist an $n \geq 0$ and an immersion
$i : X \to \mathbf{P}^n_S$ over $S$ such that
$\mathcal{L}^{\otimes d} \cong i^*\mathcal{O}_{\mathbf{P}^n_S}(1)$.
\end{lemma}

\begin{proof}
Let
$A = \Gamma_*(X, \mathcal{L}) =
\bigoplus_{d \geq 0} \Gamma(X, \mathcal{L}^{\otimes d})$.
By Properties, Proposition \ref{properties-proposition-characterize-ample}
the set of affine opens $X_a$ with $a \in A_{+}$ homogeneous forms
a basis for the topology of $X$. Hence we can find finitely
many such elements $a_0, \ldots, a_n \in A_{+}$ such that
\begin{enumerate}
\item we have $X = \bigcup_{i = 0, \ldots, n} X_{a_i}$,
\item each $X_{a_i}$ is affine, and
\item each $X_{a_i}$ maps into an affine open $V_i \subset S$.
\end{enumerate}
By Lemma \ref{lemma-locally-finite-type-characterize}
we see that the ring maps
$\mathcal{O}_S(V_i) \to \mathcal{O}_X(X_{a_i})$ are
of finite type. Hence we can find finitely many
elements $f_{ij} \in \mathcal{O}_X(X_{a_i})$, $j = 1, \ldots, n_i$
which generate $\mathcal{O}_X(X_{a_i})$ as an $\mathcal{O}_S(V_i)$-algebra.
By Properties, Lemma \ref{properties-lemma-invert-s-sections}
we may write each
$f_{ij}$ as $a_{ij}/a_i^{e_{ij}}$ for some
$a_{ij} \in A_{+}$ homogeneous. Let $N$ be a positive integer which
is a common multiple of all the degrees of the elements
$a_i$, $a_{ij}$. Consider the elements
$$
a_i^{N/\deg(a_i)}, \ a_{ij}a_i^{(N/\deg(a_i)) - e_{ij}} \in A_N.
$$
By construction these generate the invertible sheaf
$\mathcal{L}^{\otimes N}$ over $X$. Hence they give rise
to a morphism
$$
j : X \longrightarrow
\mathbf{P}_S^{m}
\quad
\text{with } m = n + \sum n_i
$$
over $S$, see Constructions, Lemma \ref{constructions-lemma-projective-space}
and Definition \ref{constructions-definition-projective-space}.
Moreover, $j^*\mathcal{O}_{\mathbf{P}_S}(1) = \mathcal{L}^{\otimes N}$.
We name the homogeneous coordinates $T_0, \ldots, T_n, T_{ij}$
instead of $T_0, \ldots, T_m$.
For $i = 0, \ldots, n$ we have $j^{-1}(D_{+}(T_i)) = X_{a_i}$.
Moreover, pulling back the element $T_{ij}/T_i$ via $j^\sharp$ we
get the element $f_{ij} \in \mathcal{O}_X(X_{a_i})$.
Hence the morphism $j$ restricted to $X_{a_i}$
gives a closed immersion of $X_{a_i}$ into the affine open
$D_{+}(T_i) \cap \mathbf{P}^m_{V_i}$ of $\mathbf{P}^N_S$.
Hence we conclude that the morphism $j$ is an immersion.
This implies the lemma holds for some $d$ and $n$ which is enough
in virtually all applications.

\medskip\noindent
This proves that for one $d_2 \geq 1$
(namely $d_2 = N$ above), some $m \geq 0$ there exists some
immersion $j : X \to \mathbf{P}^m_S$ given by global sections
$s'_0, \ldots, s'_m \in \Gamma(X, \mathcal{L}^{\otimes d_2})$.
By Properties, Proposition \ref{properties-proposition-characterize-ample}
we know there exists an integer
$d_1$ such that $\mathcal{L}^{\otimes d}$ is globally generated
for all $d \geq d_1$. Set $d_0 = d_1 + d_2$. We claim that
the lemma holds with this value of $d_0$. Namely, given
an integer $d \geq d_0$ we may choose $s''_1, \ldots, s''_t
\in \Gamma(X, \mathcal{L}^{\otimes d - d_2})$ which generate
$\mathcal{L}^{\otimes d - d_2}$ over $X$. Set $k = (m + 1)t$ and
denote $s_0, \ldots, s_k$ the collection of sections
$s'_\alpha s''_\beta$, $\alpha = 0, \ldots, m$,
$\beta = 1, \ldots, t$. These generate $\mathcal{L}^{\otimes d}$
over $X$ and therefore define a morphism
$$
i : X \longrightarrow \mathbf{P}^{k - 1}_S
$$
such that $i^*\mathcal{O}_{\mathbf{P}^n_S}(1) \cong \mathcal{L}^{\otimes d}$.
To see that $i$ is an immersion, observe that $i$ is the composition
$$
X \longrightarrow \mathbf{P}^m_S \times_S \mathbf{P}^{t - 1}_S
\longrightarrow \mathbf{P}^{k - 1}_S
$$
where the first morphism is $(j, j')$ with $j'$ given by
$s''_1, \ldots, s''_t$ and the
second morphism is the Segre embedding
(Constructions, Lemma \ref{constructions-lemma-segre-embedding}).
Since $j$ is an immersion, so is $(j, j')$
(apply Lemma \ref{lemma-immersion-permanence}
to $X \to \mathbf{P}^m_S \times_S \mathbf{P}^{t - 1}_S
\to \mathbf{P}^m_S$). Thus $i$ is a composition of
immersions and hence an immersion
(Schemes, Lemma \ref{schemes-lemma-composition-immersion}).
\end{proof}

\begin{lemma}
\label{lemma-finite-type-over-affine-ample-very-ample}
Let $f : X \to S$ be a morphism of schemes.
Let $\mathcal{L}$ be an invertible $\mathcal{O}_X$-module.
Assume $S$ affine and $f$ of finite type.
The following are equivalent
\begin{enumerate}
\item $\mathcal{L}$ is ample on $X$,
\item $\mathcal{L}$ is $f$-ample,
\item $\mathcal{L}^{\otimes d}$ is $f$-very ample for some $d \geq 1$,
\item $\mathcal{L}^{\otimes d}$ is $f$-very ample for all $d \gg 1$,
\item for some $d \geq 1$ there exist $n \geq 1$ and an immersion
$i : X \to \mathbf{P}^n_S$ such that
$\mathcal{L}^{\otimes d} \cong i^*\mathcal{O}_{\mathbf{P}^n_S}(1)$, and
\item for all $d \gg 1$ there exist $n \geq 1$ and an immersion
$i : X \to \mathbf{P}^n_S$ such that
$\mathcal{L}^{\otimes d} \cong i^*\mathcal{O}_{\mathbf{P}^n_S}(1)$.
\end{enumerate}
\end{lemma}

\begin{proof}
The equivalence of (1) and (2) is Lemma \ref{lemma-ample-over-affine}.
The implication (2) $\Rightarrow$ (6) is
Lemma \ref{lemma-quasi-projective-finite-type-over-S}.
Trivially (6) implies (5).
As $\mathbf{P}^n_S$ is a projective bundle over $S$ (see
Constructions, Lemma \ref{constructions-lemma-projective-space-bundle})
we see that
(5) implies (3) and (6) implies (4) from the definition of a
relatively very ample sheaf.
Trivially (4) implies (3). To finish we have to show that
(3) implies (2) which follows from Lemma \ref{lemma-ample-very-ample}
and Lemma \ref{lemma-ample-power-ample}.
\end{proof}

\begin{lemma}
\label{lemma-finite-type-ample-very-ample}
Let $f : X \to S$ be a morphism of schemes.
Let $\mathcal{L}$ be an invertible $\mathcal{O}_X$-module.
Assume $S$ quasi-compact and $f$ of finite type.
The following are equivalent
\begin{enumerate}
\item $\mathcal{L}$ is $f$-ample,
\item $\mathcal{L}^{\otimes d}$ is $f$-very ample for some $d \geq 1$,
\item $\mathcal{L}^{\otimes d}$ is $f$-very ample for all $d \gg 1$.
\end{enumerate}
\end{lemma}

\begin{proof}
Trivially (3) implies (2). Lemma \ref{lemma-ample-very-ample} guarantees that
(2) implies (1) since a morphism of finite type is quasi-compact
by definition. Assume that $\mathcal{L}$ is $f$-ample. Choose a finite affine
open covering $S = V_1 \cup \ldots \cup V_m$. Write $X_i = f^{-1}(V_i)$.
By Lemma \ref{lemma-finite-type-over-affine-ample-very-ample} above we see
there exists a $d_0$ such that $\mathcal{L}^{\otimes d}$ is
relatively very ample on $X_i/V_i$ for all $d \geq d_0$. Hence we conclude
(1) implies (3) by Lemma \ref{lemma-relatively-very-ample}.
\end{proof}

\noindent
The following two lemmas provide the most used and most useful
characterizations of relatively very ample and relatively ample
invertible sheaves when the morphism is of finite type.

\begin{lemma}
\label{lemma-characterize-very-ample-on-finite-type}
Let $f : X \to S$ be a morphism of schemes.
Let $\mathcal{L}$ be an invertible sheaf on $X$.
Assume $f$ is of finite type.
The following are equivalent:
\begin{enumerate}
\item $\mathcal{L}$ is $f$-relatively very ample, and
\item there exist an open covering $S = \bigcup V_j$,
for each $j$ an integer $n_j$, and immersions
$$
i_j :
X_j = f^{-1}(V_j) = V_j \times_S X
\longrightarrow
\mathbf{P}^{n_j}_{V_j}
$$
over $V_j$ such that
$\mathcal{L}|_{X_j} \cong i_j^*\mathcal{O}_{\mathbf{P}^{n_j}_{V_j}}(1)$.
\end{enumerate}
\end{lemma}

\begin{proof}
We see that (1) implies (2) by taking an affine open covering of $S$
and applying Lemma \ref{lemma-very-ample-finite-type-over-affine} to
each of the restrictions of $f$ and
$\mathcal{L}$. We see that (2) implies (1) by
Lemma \ref{lemma-relatively-very-ample}.
\end{proof}

\begin{lemma}
\label{lemma-characterize-ample-on-finite-type}
Let $f : X \to S$ be a morphism of schemes.
Let $\mathcal{L}$ be an invertible sheaf on $X$.
Assume $f$ is of finite type.
The following are equivalent:
\begin{enumerate}
\item $\mathcal{L}$ is $f$-relatively ample, and
\item there exist an open covering $S = \bigcup V_j$,
for each $j$ an integers $d_j \geq 1$,
$n_j \geq 0$, and immersions
$$
i_j :
X_j = f^{-1}(V_j) = V_j \times_S X
\longrightarrow
\mathbf{P}^{n_j}_{V_j}
$$
over $V_j$ such that
$\mathcal{L}^{\otimes d_j}|_{X_j} \cong
i_j^*\mathcal{O}_{\mathbf{P}^{n_j}_{V_j}}(1)$.
\end{enumerate}
\end{lemma}

\begin{proof}
We see that (1) implies (2) by taking an affine open covering of $S$
and applying Lemma \ref{lemma-finite-type-over-affine-ample-very-ample} to
each of the restrictions of $f$ and
$\mathcal{L}$. We see that (2) implies (1) by
Lemma \ref{lemma-characterize-relatively-ample}.
\end{proof}

\begin{lemma}
\label{lemma-invertible-add-enough-ample-very-ample}
Let $f : X \to S$ be a morphism of schemes.
Let $\mathcal{N}$, $\mathcal{L}$ be invertible $\mathcal{O}_X$-modules.
Assume $S$ is quasi-compact, $f$ is of finite type, and $\mathcal{L}$
is $f$-ample. Then
$\mathcal{N} \otimes_{\mathcal{O}_X} \mathcal{L}^{\otimes d}$
is $f$-very ample for all $d \gg 1$.
\end{lemma}

\begin{proof}
By Lemma \ref{lemma-characterize-very-ample-on-finite-type}
we reduce to the case $S$ is affine. Combining
Lemma \ref{lemma-finite-type-over-affine-ample-very-ample} and
Properties, Proposition \ref{properties-proposition-characterize-ample}
we can find an integer $d_0$ such that
$\mathcal{N} \otimes \mathcal{L}^{\otimes d_0}$
is globally generated. Choose global sections
$s_0, \ldots, s_n$ of $\mathcal{N} \otimes \mathcal{L}^{\otimes d_0}$
which generate it. This determines a morphism
$j : X \to \mathbf{P}^n_S$ over $S$. By
Lemma \ref{lemma-finite-type-over-affine-ample-very-ample}
we can also pick an integer $d_1$ such that for all $d \geq d_1$
there exist sections $t_{d, 0}, \ldots, t_{d, n(d)}$
of $\mathcal{L}^{\otimes d}$ which generate it and define
an immersion
$$
j_d = \varphi_{\mathcal{L}^{\otimes d}, t_{d, 0}, \ldots, t_{d, n(d)}} :
X
\longrightarrow
\mathbf{P}^{n(d)}_S
$$
over $S$. Then for $d \geq d_0 + d_1$ we can consider the
morphism
$$
\varphi_{\mathcal{N} \otimes \mathcal{L}^{\otimes d},
s_j \otimes t_{d - d_0, i}} :
X
\longrightarrow
\mathbf{P}^{(n + 1)(n(d - d_0) + 1) - 1}_S
$$
This morphism is an immersion as it is the composition
$$
X \to \mathbf{P}^n_S \times_S \mathbf{P}^{n(d - d_0)}_S
\to \mathbf{P}^{(n + 1)(n(d - d_0) + 1) - 1}_S
$$
where the first morphism is $(j, j_{d - d_0})$ and
the second is the Segre embedding
(Constructions, Lemma \ref{constructions-lemma-segre-embedding}).
Since $j$ is an immersion, so is $(j, j_{d - d_0})$
(apply Lemma \ref{lemma-immersion-permanence}). We have a composition of
immersions and hence an immersion
(Schemes, Lemma \ref{schemes-lemma-composition-immersion}).
\end{proof}










